\documentclass[10pt,a4paper]{amsart}
\usepackage[margin=19mm]{geometry}
\usepackage{amsmath,amssymb,mathtools}
\usepackage[scr=rsfs,cal=euler]{mathalfa}
\usepackage{xfrac}
\usepackage[unicode,hidelinks,bookmarks=false]{hyperref}
\newcommand{\ie}{i.e.\ }
\newcommand{\cf}{cf.\ }
\newcommand{\resp}{respectively\ }
\newcommand{\Subsection}{\subsection}
\providecommand{\nobreakdash}{\nobreak\hbox{-}}
\setlength{\emergencystretch}{2em}
\multlinegap0pt
\usepackage{bm}
\usepackage{algorithm}\usepackage{algpseudocode}
\usepackage{tikz}
\usepackage[normalem]{ulem}
\usepackage{amssymb}
\usepackage{xcolor}
\usepackage[shortlabels]{enumitem}
\usepackage{csquotes}
\newtheorem{assumption}{Assumption}
\newtheorem{lemma}{Lemma}
\usepackage{cleveref}
\crefname{assumption}{Assumption}{Assumptions}
\crefname{lemma}{Lemma}{Lemmas}
\newcommand{\Cbb}{\mathbb{C}}
\newcommand{\History}{\mathcal{H}}
\newcommand{\Ibb}{\mathbb{I}}
\newcommand{\Pbb}{\mathbb{P}}
\newcommand{\Tbb}{\mathbb{T}}
\newcommand{\constant}{\Cbb}
\newcommand{\deltam}{\lVert \Delta_m \rVert}
\newcommand{\expec}{\mathbb{E}}
\newcommand{\history}{\hbar}
\newcommand{\indexeddata}{\left\{(X_0,a_0),\dots,(X_m,a_m)\right\}}
\newcommand{\lc}{\left\{}
\newcommand{\rc}{\right\}}
\newcommand{\lp}{\left(}
\newcommand\numberthis{\addtocounter{equation}{1}\tag{\theequation}}
\newcommand{\pow}[1]{^{(#1)}}
\newcommand{\prob}{\Pbb}
\newcommand{\probl}{\prob}
\newcommand{\rp}{\right)}
\title{Offline Estimation of Controlled Markov Chains: Minimaxity and Sample Complexity}
\author{Imon Banerjee, Harsha Honnappa, Vinayak Rao}
\date{Source: arXiv:2211.07092v5 (CC BY 4.0)}
\begin{document}
\maketitle

\noindent\emph{Source and licence.} Offline Estimation of Controlled Markov Chains: Minimaxity and Sample Complexity. Version arXiv:2211.07092v5 (CC BY 4.0), distributed by arXiv under the Creative Commons Attribution 4.0 International licence, http://creativecommons.org/licenses/by/4.0/, as recorded in the arXiv licence metadata for this exact version and shown on the abstract page https://arxiv.org/abs/2211.07092v5. Full licence notice: \texttt{licenses/arxiv-2211.07092-CC-BY-4.0.txt}.\par\smallskip

\noindent\textit{Editorial scope.} Counted unit: the article's lemma lemma:kontorovich (source0.tex lines 1571--1578). The article's complete original proof of it is delivered with it (source0.tex lines 3218--3316, one \texttt{proof} environment of 99 lines, which the article places away from the statement). No mathematics is written, repaired or re-derived: the statement and the proof are the article's own lines, in the article's own notation, and no retained line is substituted or re-flowed.  Retained uncounted as the source-local context the proof needs: lines 719--721; lines 1113--1118; lines 1134--1138; lines 1636--1642.  Nothing was omitted from any retained span, so the package carries no omission record.  No retained line contains a citation, so the wrapper carries no bibliography and no bibliography entry.  No retained line contains a figure environment, an image-inclusion command or drawing code, so no image is delivered and the package declares no asset dependency.  Editorial formatting only, in addition: an AMS \texttt{amsart} preamble that restates the article's own declarations and macros used by the retained lines, namely the environments \texttt{assumption}, \texttt{lemma} on separate counters, together with the standard packages those lines need.  The retained environments print in this document as Assumption 1, Lemma 1 in order of appearance, whereas the article prints the same items with the numbering of its own full document.  The complete native source of the article as distributed in the arXiv e-print for this exact version is bundled unchanged as \texttt{sources/133431/source0.tex}.
    \begin{align*}~\label{def:weak-mixing}
    \Bar{\eta}_{i,j}:= \underset{\color{black}\substack{\mathbb{T},s_1,s_2,l_1,l_2,\history_0^{i-1},\\\prob\lp X_i=s_1,a_i=l_1,\History_0^{i-1}=\history_0^{i-1}\rp>0,\\\prob\lp  X_i=s_2,a_i=l_2,\History_0^{i-1}=\history_0^{i-1}\rp>0}}{\sup} \eta_{i,j}(\mathbb{T},s_1,s_2,l_1,l_2,\history_0^{i-1}).\numberthis
\end{align*}

\begin{assumption}[Mixing of controls]~\label{assume:control-mixing} There exists a constant $\constant\geq 0$ such that
\begin{align*}
    \sup_{1\leq i\leq{\color{black} \infty}}\sum_{j=1}^{{\color{black} \infty}} \sum_{p=i+j+1}^{{\color{black} \infty}} \gamma_{p,j,i}\leq \frac{\constant}{2}.
\end{align*}

\end{assumption}

\begin{assumption}[Mixing of States]~\label{assume:chain-mix} There exists a constant $\constant_{\theta}\geq 0$ such that
\[
\sup_{1\leq i\leq { \infty}}\sum_{j=i+1}^{{ \infty}} \bar \theta_{i,j}\leq\constant_{\theta}. 
\]
\end{assumption}

\begin{lemma}~\label{lemma:tv-bound}
Let $X,Y,$ and $Z$ be 3 discrete random variables in the sample space $\Omega$. Then
    \begin{align*}
        \lVert \probl(X,Y|Z=z_1)& -\probl(X,Y|Z=z_2) \rVert_{TV} \leq \lVert \probl(Y|z=z_1)-\probl(Y|Z=z_2) \rVert_{TV}\\
        & \qquad +\sup_y\lVert \probl(X|Y=y,Z=z_1)-\probl(X|Y=y,Z=z_2) \rVert_{TV}.
    \end{align*}
\end{lemma}

\begin{lemma}~\label{lemma:kontorovich}
    Let $\indexeddata$ be a sequence of stochastic random variables on a finite state space $\chi\times\Ibb$. 
    Then, for any $t > 0$ we have,
  \begin{align}
      \prob(|N_s\pow{l}-\expec[N_s\pow{l}]|>t)\leq 2\exp\left(-\frac{t^2}{2m\|\Delta_m\|^2}\right),
  \end{align}
where $\|\Delta_m\|=\underset{1\leq i\leq m}{\max} (1+ \Bar{\eta}_{i,i+1}+ \Bar{\eta}_{i,i+2}+\dots \Bar{\eta}_{i,m})$, and $\Bar{\eta}_{i,j}$ is as defined in \cref{def:weak-mixing}.
\end{lemma}

\begin{proof} \
Recall from Lemma \ref{lemma:kontorovich} the definition of $\deltam$.
\[
\|\Delta_m\|:=\underset{1\leq i\leq m}{\max} (1+ \Bar{\eta}_{i,i+1}+ \Bar{\eta}_{i,i+2}+\dots \Bar{\eta}_{i,m}),
\]
where, 
\begin{align*}
    \Bar{\eta}_{i,j}:= \underset{\mathbb{T},s_1,s_2,l_1,l_2,\history_0^{i-1}}{\sup} \eta_{i,j},
\end{align*}
and, 
\begin{align*}
    \eta_{i,j} & := \left|\prob\left((X_m,a_m,\dots, X_j,a_j)\in \mathbb{T}|X_i=s_1,a_i=l_1,\History_0^{i-1}=\history_0^{i-1}\right)\right.\\
    & \left.\qquad \qquad -\prob\left((X_m,a_m,\dots,X_j,a_j)\in\mathbb{T}|X_i=s_2,a_i=l_2,\History_0^{i-1}=\history_0^{i-1}\right)\right|.
\end{align*}
For any fixed $s_1,s_2\in \chi$, $a_1,a_2\in\Ibb$, and $\history_0^{i-1}\in (\chi\times\Ibb)^{i}$, define the total variation distance $\eta_{i,j}\pow{m}$ as, 
\begin{align*}
\eta_{i,j}\pow{m} & := \|\probl\left((X_m,a_m,\dots, X_j,a_j)|X_i=s_1,a_i=l_1,\History_1^{i-1}=\history_0^{i-1}\right)\\
    &  \qquad -\probl\left((X_m,a_m,\dots,X_j,a_j)|X_i=s_2,a_i=l_2,\History_1^{i-1}=\history_0^{i-1}\right)\|_{TV}.
\end{align*}
We note that for each fixed $m$, $\bar \eta_{i,j}$ can be recovered from the corresponding $\eta_{i,j}\pow{m}$ as 
\begin{align*}
    \bar \eta_{i,j} = \sup_{s_1,s_2,l_1,l_2,\history_0^{i-1}} \eta_{i,j}\pow{m}. \numberthis\label{eq:db-eq4}
\end{align*}
Thus is is enough to find an upper bound for $\eta_{i,j}\pow{m}$. 
For the sake of simplicity, consider the case when $m=2$. Then, 
\begin{align*}
    \eta_{1,2}\pow{2} & =\bigg\|\probl\lp(X_2,a_2)|X_1=s_1,a_1=l_1,X_0=s_0,a_0=l_0\rp\\
    & \qquad -\probl\lp(X_2,a_2)|X_1=s_2,a_1=l_2,X_0=s_0,a_0=l_0\rp\bigg\|_{TV}.
\end{align*}
Use Lemma \ref{lemma:tv-bound} by substituting $X=a_2$, $Y=X_2$ and $Z=(X_1,a_1,X_0,a_0)$ with $z_1$ and $z_2$ defined accordingly.
This reduces the $\eta$ mixing coefficient into a sum of the $\gamma$ mixing coefficients and $\theta$ mixing coefficients. 
To be precise, 
\begin{align*}
    \eta_{1,2}\pow{2}& \leq \sup_{s''}\| \probl\lp a_2|X_2=s'',\numberthis\label{eq:db-eq1} X_1=s_1,a_1=l_1,X_0=s_0,a_0=l_0\rp\\
    & \qquad -\probl\lp a_2|X_2=s'', X_1=s_2,a_1=l_2,X_0=s_0,a_0=l_0\rp\|_{TV}\\
    & + \|\probl\lp X_2\in \Tbb_2| X_1=s_1,a_1=l_1,X_0=s_0,a_0=l_0\rp\\
    & \qquad -\probl\lp X_2\in \Tbb_2| X_1=s_2,a_1=l_2,X_0=s_0,a_0=l_0\rp\|_{TV}.
\end{align*}
We can decompose the first term in \cref{eq:db-eq1} as,
\begin{align*}
    &\sup_{s''}\| \probl\lp a_2|X_2=s'', X_1=s_1,a_1=l_1,X_0=s_0,a_0=l_0\rp\\
    & \qquad -\probl\lp a_2|X_2=s'', X_1=s_2,a_1=l_2,X_0=s_0,a_0=l_0\rp\|_{TV}\\
    & \ \leq   \sup_{s''}\| \probl\lp a_2|X_2=s'' X_1=s_1,a_1=l_1,X_0=s_0,a_0=l_0\rp- \probl\lp a_2|X_2=s''\rp\|_{TV}\\
    & \qquad + \sup_{s''}\|\probl\lp a_2|X_2=s'', X_1=s_2,a_1=l_2,X_0=s_0,a_0=l_0\rp-\probl\lp a_2|X_2=s''\rp\|_{TV}\\
    & \ \leq 2\gamma_{2,1,1}\numberthis\label{eq:db-eq3},
\end{align*}
where the last step follows from taking supremum over $s_1,l_1,s_2,l_2,s_0$,and $a_0$ on both sides of the inequality. 
It follows that, when $m=2$,
\[
 \bar\eta_{1,2}\leq \gamma_{2,1,1}+\bar\theta_{1,2}.
\]
We now establish the following recursion,
\begin{align*}
 \eta_{i,j}^{(m+1)} \leq \eta_{i,j}^{(m)}+2\gamma_{m+1,j,i}.\numberthis\label{eq:db-eq5}    
\end{align*}

Fix an arbitrary $m\geq 3$ and fix $i,j\in{1,2,\dots,m}$ such that $i<j$ and
decompose $\eta_{i,j}\pow{m+1}$ using Lemma \ref{lemma:tv-bound} by taking $X=(X_{m+1},a_{m+1})$, and defining $Y$ and $Z$ accordingly to get,
\begin{align*}
    \eta_{i,j}^{(m+1)}& \leq \eta_{i,j}^{(m)}\\
    & \qquad + \|\probl(X_{m+1},a_{m+1}|X_{m}=x_{m},a_{m}=l_{m},X_i=s_1,a_i=l_1,\dots)\\
    & \qquad \qquad -\probl(X_{m+1},a_{m+1}|X_{m}=x_{m},a_{m}=l_{m},X_i=s_2,a_i=l_2,\dots\|_{TV}.
\end{align*}
where we have replaced the terms common to both the conditionals by $\dots$ for convenience of notation.

Using Lemma \ref{lemma:tv-bound} again on the second term by taking $X=a_{m+1}$, $Y=X_{m+1}$, and $Z$ as all the random variables in the conditional, we get the following upper bound,
\begin{align*}
    \eta_{i,j}^{(m+1)}& \leq \eta_{i,j}^{(m)} + \sup_{x_{m+1}\in \chi}\|\probl(a_{m+1}|X_{m+1}=x_{m+1},X_{m}=x_{m},a_{m}=l_{m},X_i=s_1,a_i=l_1,\dots)\\
    & \qquad \qquad -\probl(a_{m+1}|X_{m+1}=x_{m+1},X_{m}=x_{m},a_{m}=l_{m},X_i=s_2,a_i=l_2,\dots\|_{TV}\\
    & \qquad + \|\probl(X_{m+1}|X_{m}=x_{m},a_{m}=l_{m},X_i=s_1,a_i=l_1,\dots)\\
    & \qquad \qquad -\probl(X_{m+1}|X_{m}=x_{m},a_{m}=l_{m},X_i=s_2,a_i=l_2,\dots\|_{TV}.\numberthis\label{eq:db-eq2}
\end{align*}
Since $X_{m+1}$ is Markovian conditional on $X_{m}$ and $a_{m}$, the third term vanishes. Proceeding similarly to \cref{eq:db-eq3}, the second term can be upper bounded by $\gamma_{m+1,j,i}$.
This proves the recursion in \cref{eq:db-eq5}.

Applying the recursion multiple times we are left with, 
\[
\eta_{i,j}^{(m)}\leq \sum_{p=i+j+1}^m 2\gamma_{p,j,i}+\eta_{i,j}\pow{j}.
\]
Proceeding similarly to \cref{eq:db-eq1}, we can decompose $\eta_{i,j}\pow{j}$ to get, 
\[
\eta_{i,j}^{(j)}\leq 2\gamma_{j,i,i}+\bar\theta_{i,j}.
\]
The terms in the right hand side of the previous equation is a constant. 
Following the arguments of \cref{eq:db-eq4}, we can now take an appropriate supremum on the left hand side to recover $\bar \eta_{i,j}$. 
Therefore, 
\begin{align*}
     \bar\eta_{i,j}\leq 2\sum_{p=i+j}^m\gamma_{p,j,i}+\bar\theta_{i,j}.\numberthis\label{eq:lem-dbeq1}   
\end{align*}
By substituting the upper bound into the expression of $\deltam$, it follows that, 
\[
    \deltam \leq \sup_{1\leq i\leq m}\lc 1+2\sum_{j>i}\sum_{p=i+j}^m\gamma_{p,j,i}+\sum_{j>i}\bar\theta_{i,j}\rc,
\]
From Assumptions \ref{assume:control-mixing} and \ref{assume:chain-mix} it now follows that,
\[
     \deltam \leq 1+\constant+\constant_{\theta}, 
\]
which completes the proof.
\end{proof}

\end{document}
